\documentclass{article}
\usepackage{amssymb}
\usepackage{algorithm}
\usepackage{algorithmic}
\begin{document}

\begin{algorithm}
\caption{Euclidean Algorithm}\label{alg:euclid}
\begin{algorithmic}[1]
\REQUIRE $a, b \in \mathbb{N}$ with $a \ge b$
\ENSURE $\gcd(a, b)$
\STATE $r \leftarrow a \bmod b$\label{line:alc:mod}
\WHILE{$r \ne 0$}\label{line:alc:while}
  \STATE $a \leftarrow b$
  \STATE $b \leftarrow r$
  \STATE $r \leftarrow a \bmod b$
\ENDWHILE\label{line:alc:endwhile}
\RETURN $b$\label{line:alc:ret}
\end{algorithmic}
\end{algorithm}

\begin{algorithmic}[1]
\STATE Standalone computation\label{line:alc:standalone}
\end{algorithmic}

In Algorithm~\ref{alg:euclid}, line~\ref{line:alc:mod} initializes the remainder.
The loop spanning lines~\ref{line:alc:while}--\ref{line:alc:endwhile} reduces the inputs,
and line~\ref{line:alc:ret} returns the result.
Standalone line~\ref{line:alc:standalone} executes outside any float.

\end{document}
